\documentclass[12pt,a4paper]{article}

% ---- Packages ----------------------------------------------------------------------------------------------------------------------------------
\usepackage[margin=2.5cm]{geometry}
\usepackage{amsmath, amssymb}
\usepackage{graphicx}
\usepackage{hyperref}
\usepackage{xcolor}
\usepackage{listings}
\usepackage{booktabs}
\usepackage{array}
\usepackage{longtable}
\usepackage{enumitem}
\usepackage{fancyhdr}
\usepackage{titlesec}
\usepackage{mdframed}
\usepackage{fontenc}
\usepackage{inputenc}

% ---- Colors --------------------------------------------------------------------------------------------------------------------------------------
\definecolor{angkorblue}{RGB}{0,70,127}
\definecolor{angkorgold}{RGB}{180,140,0}
\definecolor{codegreen}{RGB}{0,120,0}
\definecolor{codegray}{RGB}{100,100,100}
\definecolor{codeblue}{RGB}{0,0,180}
\definecolor{lightgray}{RGB}{245,245,245}

% ---- Code Listing Style ----------------------------------------------------------------------------------------------------------------
\lstdefinestyle{pythonstyle}{
    backgroundcolor=\color{lightgray},
    commentstyle=\color{codegreen},
    keywordstyle=\color{codeblue}\bfseries,
    stringstyle=\color{red},
    basicstyle=\ttfamily\small,
    breakatwhitespace=false,
    breaklines=true,
    captionpos=b,
    keepspaces=true,
    numbers=left,
    numberstyle=\tiny\color{codegray},
    showspaces=false,
    showstringspaces=false,
    showtabs=false,
    tabsize=4,
    language=Python,
    frame=single,
    rulecolor=\color{angkorblue}
}
\lstset{style=pythonstyle}

% ---- Header/Footer --------------------------------------------------------------------------------------------------------------------------
\pagestyle{fancy}
\fancyhf{}
\rhead{\textcolor{angkorblue}{ANGKOR Project Documentation}}
\lhead{\textcolor{angkorblue}{\leftmark}}
\rfoot{\thepage}
\lfoot{\textcolor{codegray}{Institute of Technology of Cambodia}}

% ---- Section Colors ------------------------------------------------------------------------------------------------------------------------
\titleformat{\section}
  {\Large\bfseries\color{angkorblue}}
  {\thesection}{1em}{}[\titlerule]

\titleformat{\subsection}
  {\large\bfseries\color{angkorblue!80}}
  {\thesubsection}{1em}{}

\titleformat{\subsubsection}
  {\normalsize\bfseries\color{angkorblue!60}}
  {\thesubsubsection}{1em}{}

% ---- Hyperref Setup ------------------------------------------------------------------------------------------------------------------------
\hypersetup{
    colorlinks=true,
    linkcolor=angkorblue,
    urlcolor=angkorgold,
    citecolor=angkorblue
}

% ----------------------------------------------------------------------------------------------------------------------------------------------------------
\begin{document}

% ---- Title Page --------------------------------------------------------------------------------------------------------------------------------
\begin{titlepage}
    \centering
    \vspace*{2cm}

    {\Huge\bfseries\color{angkorblue} ANGKOR}\\[0.5cm]
    {\Large\color{angkorgold}
        Advanced Neutron Group-diffusion\\
        K-eigenvalue Of Reactors}\\[1.5cm]

    \rule{\linewidth}{0.5mm}\\[0.4cm]
    {\large\bfseries Project Development Summary \& Documentation}\\
    {\normalsize Technical Reference, Decisions, and Future Roadmap}\\
    \rule{\linewidth}{0.5mm}\\[2cm]

    \begin{minipage}{0.6\textwidth}
        \centering
        {\large\textbf{MUTH Boravy}}\\
        Nuclear Engineering Student\\
        Institute of Technology of Cambodia (ITC)\\
        Phnom Penh, Cambodia\\[0.5cm]
        \texttt{\href{https://github.com/muthboravy007/ANGKOR}
        {github.com/muthboravy007/ANGKOR}}\\[0.5cm]
        \textit{Developed with AI assistance (Claude, Anthropic)}
    \end{minipage}

    \vfill
    {\large March 2026}
\end{titlepage}

% ---- Table of Contents ------------------------------------------------------------------------------------------------------------------
\tableofcontents
\newpage

% =============================================================================
\section{Project Overview}
% =============================================================================

ANGKOR is an open-source deterministic nuclear reactor physics code developed
at the Institute of Technology of Cambodia (ITC). The name honors
\textbf{Angkor Wat} --- the greatest symbol of Khmer civilization --- as a
statement that Cambodia's scientific legacy continues in the modern era.

\begin{mdframed}[linecolor=angkorblue, linewidth=2pt]
\textbf{Mission Statement:} To place Cambodia on the world map of nuclear
reactor physics analysis, providing an open-source alternative to commercial
codes such as CASMO and SIMULATE, freely accessible to researchers,
students, and engineers worldwide.
\end{mdframed}

\subsection{Code Identity}

\begin{center}
\begin{tabular}{ll}
\toprule
\textbf{Full Name} & Advanced Neutron Group-diffusion K-eigenvalue Of Reactors \\
\textbf{Language} & Python 3.13 (NumPy, SciPy, Matplotlib) \\
\textbf{Method} & Deterministic: Finite Difference + Power Iteration \\
\textbf{License} & MIT (free to use, modify, distribute) \\
\textbf{Repository} & \url{https://github.com/muthboravy007/ANGKOR} \\
\textbf{Institution} & Institute of Technology of Cambodia (ITC) \\
\textbf{Country} & Cambodia \\ 
\bottomrule
\end{tabular}
\end{center}

% =============================================================================
\section{Development History}
% =============================================================================

The project evolved from a simple 1D slab solver (originally named PyReactor)
into a multi-group 2D diffusion code system. The development followed a
structured pedagogical approach: build, validate, extend.

\subsection{Version Timeline}

\begin{longtable}{p{2.5cm} p{4cm} p{7cm}}
\toprule
\textbf{Version} & \textbf{Milestone} & \textbf{Description} \\
\midrule
v1.0 & 1D Foundation & 1D slab geometry, 2-group diffusion, power iteration.
Validated: 0 pcm vs analytical solution. \\[6pt]
v1.1 & 2D Geometry & 2D rectangular geometry engine with material map,
interactive visualizer (color-coded, zoom/pan). \\[6pt]
v1.2 & YAML Input & Human-readable YAML input system. Auto-detects 2-group
vs multi-group format. \\[6pt]
v1.3 & 2D Solver & 2D finite difference solver with reflective and vacuum
boundary conditions. IAEA benchmark run. \\[6pt]
v1.4 & Output Module & Professional flux maps, power distribution plots,
centerline profiles, text reports. \\[6pt]
v1.5 & Unit Tests & pytest suite: geometry, materials, solver tests. \\[6pt]
v2.0 & Multi-Group & G-group solver (SolverMG) supporting arbitrary
number of energy groups. \\[6pt]
v2.1 & XS Reader & OpenMC HDF5 nuclear data reader with group condensation
using flux-weighted averages. \\[6pt]
v2.2 & CMFD Start & Coarse Mesh Finite Difference acceleration framework
initiated (in progress). \\
\bottomrule
\end{longtable}

% =============================================================================
\section{Physics Foundation}
% =============================================================================

\subsection{Governing Equation}

ANGKOR solves the multi-group neutron diffusion equation in 2D Cartesian
geometry. For group $g$ ($g = 1, 2, \ldots, G$):

\begin{equation}
    -\nabla \cdot D_g(\mathbf{r}) \nabla \phi_g(\mathbf{r})
    + \Sigma_{r,g}(\mathbf{r}) \phi_g(\mathbf{r})
    = \chi_g \frac{1}{k} \sum_{g'=1}^{G} \nu\Sigma_{f,g'}(\mathbf{r})
      \phi_{g'}(\mathbf{r})
    + \sum_{\substack{g'=1\\g'\neq g}}^{G}
      \Sigma_{s,g'\to g}(\mathbf{r}) \phi_{g'}(\mathbf{r})
\end{equation}

where the removal cross section is:
\begin{equation}
    \Sigma_{r,g} = \Sigma_{a,g} + \sum_{g'\neq g} \Sigma_{s,g\to g'}
\end{equation}

\subsection{Finite Difference Discretization}

The Laplacian is approximated using the 5-point stencil on a uniform mesh
with spacing $\Delta x$, $\Delta y$:

\begin{equation}
    \nabla^2 \phi \approx
    \frac{\phi_{i+1,j} - 2\phi_{i,j} + \phi_{i-1,j}}{\Delta x^2}
    + \frac{\phi_{i,j+1} - 2\phi_{i,j} + \phi_{i,j-1}}{\Delta y^2}
\end{equation}

Each cell $(i,j)$ is connected to its four neighbors (west, east, south,
north). The center coefficient for group $g$ at cell $n$ is:

\begin{equation}
    C_g = \frac{2D_g}{\Delta x^2} + \frac{2D_g}{\Delta y^2}
          + \Sigma_{a,g} + \sum_{g'>g} \Sigma_{s,g\to g'}
\end{equation}

\subsection{Boundary Conditions}

\begin{center}
\begin{tabular}{lll}
\toprule
\textbf{Type} & \textbf{Mathematical Form} & \textbf{FDM Implementation} \\
\midrule
Vacuum (zero flux) & $\phi = 0$ at boundary & Add $D/\Delta x^2$ to diagonal \\
Reflective (zero current) & $\partial\phi/\partial n = 0$ & No leakage term added \\
Extrapolated & $\phi = 0$ at $d_{ex} = 2.1312D$ beyond & $D/(\Delta x(\Delta x/2 + d_{ex}))$ \\
\bottomrule
\end{tabular}
\end{center}

\subsection{Power Iteration Algorithm}

The eigenvalue equation $\mathbf{A}\boldsymbol{\phi} =
\frac{1}{k}\mathbf{F}\boldsymbol{\phi}$ is solved iteratively:

\begin{enumerate}
    \item Initialize: $k^{(0)} = 1.0$, $\boldsymbol{\phi}^{(0)} = \mathbf{1}$
    \item Compute fission source: $\mathbf{S}^{(n)} = \mathbf{F}\boldsymbol{\phi}^{(n)}$
    \item Solve linear system: $\mathbf{A}\boldsymbol{\phi}^{(n+1)} = \mathbf{S}^{(n)}/k^{(n)}$
    \item Update eigenvalue:
        $k^{(n+1)} = k^{(n)} \dfrac{\sum \mathbf{F}\boldsymbol{\phi}^{(n+1)}}
        {\sum \mathbf{F}\boldsymbol{\phi}^{(n)}}$
    \item Check convergence: $|k^{(n+1)} - k^{(n)}| < \varepsilon$
    \item Normalize flux and repeat until converged
\end{enumerate}

% =============================================================================
\section{Code Architecture}
% =============================================================================

\subsection{Module Structure}

\begin{verbatim}
ANGKOR/
|-- angkor/
|   |-- geometry_2d.py     # Region, GeometryEngine, GeometryVisualizer
|   |-- input_reader.py    # YAML parser, SolverSettings, auto-detect XS format
|   |-- solver_2d.py       # 2-group FD solver with BC support
|   |-- solver_mg.py       # G-group generalized solver
|   |-- output_2d.py       # Flux maps, power distribution, reports
|   |-- output_mg.py       # Multi-group output module
|   |-- cmfd.py            # CMFD acceleration (in progress)
|   |-- xs_reader.py       # HDF5 nuclear data reader
|   |-- geometry.py        # 1D geometry (legacy)
|   `-- solver_1d.py       # 1D solver (legacy)
|-- input/
|   |-- iaea_2d.yaml       # IAEA 2D PWR benchmark (validated)
|   `-- pwr_2d.yaml        # Simple PWR 2D slab test
|-- tests/
|   |-- test_geometry_2d.py
|   |-- test_solver_2d.py
|   `-- ...
|-- benchmarks/
|-- docs/
|-- main.py                # Entry point
`-- requirements.txt
\end{verbatim}

\subsection{Data Flow}

\begin{center}
\begin{tabular}{ccccc}
\texttt{YAML file}
& $\longrightarrow$ & \texttt{InputReader}
& $\longrightarrow$ & \texttt{GeometryEngine} \\
& & & & $\downarrow$ \\
\texttt{Output files}
& $\longleftarrow$ & \texttt{Output2D/MG}
& $\longleftarrow$ & \texttt{Solver2D/MG} \\
\end{tabular}
\end{center}

\subsection{Key Design Decisions}

\begin{enumerate}
    \item \textbf{Cell-centered mesh (Option 2)}: Cell centers never land on
    material boundaries, eliminating the overlap ambiguity problem. This is
    identical to the approach used in CASMO, SIMULATE, and PARCS.

    \item \textbf{Sparse matrix storage}: For a $170\times170$ mesh with
    2 groups, the full matrix would require 26.7 GB. Sparse storage reduces
    this to approximately 2.8 MB.

    \item \textbf{Auto-detection of XS format}: \texttt{InputReader}
    automatically detects 2-group (D1, D2, sigma\_a1...) vs multi-group
    (groups, D, sigma\_a...) format by checking for the \texttt{groups} key.

    \item \textbf{Separate solver files}: \texttt{solver\_2d.py} (optimized
    2-group) and \texttt{solver\_mg.py} (general G-group) coexist. The
    2-group solver preserves backward compatibility and validated results.

    \item \textbf{Reflective boundary conditions}: Implemented by omitting the
    leakage term at symmetry planes, enabling quarter-core modeling.
\end{enumerate}

% =============================================================================
\section{Validation Results}
% =============================================================================

\subsection{1D Analytical Benchmark}

\begin{center}
\begin{tabular}{lrrr}
\toprule
\textbf{Case} & \textbf{ANGKOR} & \textbf{Reference} & \textbf{Error (pcm)} \\
\midrule
Bare slab (analytical) & 1.305447 & 1.305446 & 0 \\
Moderated slab (MCNP) & 1.30866 & 1.32811 & 1464 \\
\bottomrule
\end{tabular}
\end{center}

The 1464 pcm difference with MCNP for the moderated slab is expected: it
reflects the fundamental difference between diffusion theory and Monte Carlo
transport, not a code error.

\subsection{IAEA 2D PWR Quarter-Core Benchmark}

The IAEA 2D PWR benchmark is a standard international test problem for
reactor diffusion codes. It models a quarter-core ($170 \times 170$ cm) with
three fuel types and a water reflector.

\begin{center}
\begin{tabular}{lrr}
\toprule
\textbf{Parameter} & \textbf{Value} & \textbf{Notes} \\
\midrule
Domain & $170 \times 170$ cm & Quarter-core with symmetry \\
Mesh & $170 \times 170$ & $\Delta x = \Delta y = 1.0$ cm \\
Energy groups & 2 & Fast + thermal \\
Boundary (left, bottom) & Reflective & Symmetry planes \\
Boundary (right, top) & Vacuum & Outer edges \\
Convergence tolerance & $10^{-6}$ & Both $k$ and $\phi$ \\
Iterations to converge & 73 & Power iteration \\
Runtime & $\approx 43$ seconds & Python/SciPy \\
\midrule
\textbf{ANGKOR k-eff} & \textbf{0.99837} & \\
\textbf{Reference k-eff} & \textbf{1.02959} & IAEA published value \\
\textbf{Error} & $\mathbf{-3086}$ \textbf{pcm} & \\
\bottomrule
\end{tabular}
\end{center}

\subsubsection{Error Analysis}

The $-3086$ pcm discrepancy is physically expected and well understood:

\begin{center}
\begin{tabular}{lr}
\toprule
\textbf{Error Source} & \textbf{Estimated Contribution} \\
\midrule
Diffusion approximation (vs transport) & $\sim$2000 pcm \\
2-group energy structure & $\sim$1000 pcm \\
Numerical (mesh, boundary) & $<$100 pcm \\
\midrule
\textbf{Total} & $\sim$3086 pcm \\
\bottomrule
\end{tabular}
\end{center}

Mesh convergence was verified: refining from $170\times170$ to $340\times340$
changed k-eff by only 1.1 pcm, confirming the solution is mesh-converged and
the remaining error is physics-related, not numerical.

% =============================================================================
\section{Cross Section Generation}
% =============================================================================

\subsection{OpenMC Integration}

An OpenMC-based cross section generation workflow was developed using Docker:

\begin{enumerate}
    \item Define PWR pin cell geometry (fuel pin radius = 0.4096 cm,
          pitch = 1.26 cm)
    \item Define UO2 fuel (3.1\% enriched) and water materials
    \item Run OpenMC Monte Carlo simulation (150 batches, 50 inactive,
          10,000 particles/batch)
    \item Extract flux-weighted few-group cross sections via tallies
    \item Export to YAML for ANGKOR input
\end{enumerate}

\subsection{Group Condensation Theory}

Cross sections are averaged over energy groups using flux weighting:

\begin{equation}
    \bar{\sigma}_g = \frac{\displaystyle\int_{E_{g+1}}^{E_g}
    \sigma(E)\,\phi(E)\,dE}
    {\displaystyle\int_{E_{g+1}}^{E_g} \phi(E)\,dE}
\end{equation}

The diffusion coefficient uses the harmonic mean at interfaces to ensure
current continuity:

\begin{equation}
    D_{\text{interface}} = \frac{2 D_i D_{i+1}}{D_i + D_{i+1}}
\end{equation}

\subsection{4-Group Energy Structure}

\begin{center}
\begin{tabular}{lll}
\toprule
\textbf{Group} & \textbf{Energy Range} & \textbf{Physics Region} \\
\midrule
1 & $> 0.821$ MeV & Fast (fission spectrum) \\
2 & $5.53$ keV -- $0.821$ MeV & Epithermal (slowing down) \\
3 & $0.625$ eV -- $5.53$ keV & Resonance \\
4 & $< 0.625$ eV & Thermal (Maxwellian) \\
\bottomrule
\end{tabular}
\end{center}

\subsection{Lesson Learned: Two-Step Method}

Pin cell cross sections from an infinite lattice calculation cannot be used
directly in core diffusion without a critical spectrum correction (B$^2$
correction). Proper few-group XS generation requires:

\begin{enumerate}
    \item Full assembly calculation (not just pin cell)
    \item B$^2$ search for critical spectrum
    \item Spatial homogenization (SPH factors)
\end{enumerate}

This is what CASMO does. Future work will implement this properly.

% =============================================================================
\section{CMFD Acceleration (In Progress)}
% =============================================================================

\subsection{Motivation}

Current power iteration requires 73 iterations for the IAEA benchmark.
CMFD (Coarse Mesh Finite Difference) reduces this to $\sim$10--15 iterations
by introducing a two-level solution strategy.

\subsection{Algorithm}

\begin{enumerate}
    \item Perform fine mesh sweep (Gauss-Seidel)
    \item Compute interface currents $J_{\text{fine}}$
    \item Calculate $\hat{d}$ correction coefficients:
        \begin{equation}
            \hat{d} = \frac{J_{\text{fine}} -
            \left(-\tilde{D}\,\Delta\Phi/\Delta x\right)}{\bar{\Phi}}
        \end{equation}
    \item Solve coarse mesh eigenvalue problem (small matrix)
    \item Rebalance fine mesh flux using ratio:
        \begin{equation}
            \phi_{\text{fine}}^{\text{new}} = \phi_{\text{fine}}
            \times \frac{\phi_{\text{coarse}}^{\text{new}}}
            {\phi_{\text{coarse}}^{\text{old}}}
        \end{equation}
    \item Check convergence
\end{enumerate}

\subsection{Implementation Status}

\begin{center}
\begin{tabular}{ll}
\toprule
\textbf{Component} & \textbf{Status} \\
\midrule
\texttt{\_\_init\_\_()} & Complete [done] \\
\texttt{homogenize\_flux()} & Complete [done] \\
\texttt{compute\_dhat()} & In progress \\
\texttt{build\_coarse\_matrices()} & Planned \\
\texttt{solve\_coarse()} & Planned \\
\texttt{rebalance\_flux()} & Planned \\
\texttt{accelerate()} & Planned \\
\bottomrule
\end{tabular}
\end{center}

\subsection{Expected Performance}

\begin{center}
\begin{tabular}{lrr}
\toprule
\textbf{Configuration} & \textbf{Iterations} & \textbf{Runtime} \\
\midrule
Without CMFD (current) & $\sim$73 & $\sim$43 seconds \\
With CMFD (expected) & $\sim$15 & $\sim$8 seconds \\
Speedup & $\sim$5$\times$ & \\
\bottomrule
\end{tabular}
\end{center}

% =============================================================================
\section{Teaching Methodology and Learning Outcomes}
% =============================================================================

\subsection{Teaching Approach}

The project was developed using a \textbf{fill-in-the-blank} methodology:
\begin{itemize}
    \item Instructor provides skeleton code with blanks
    \item Student fills blanks based on physical understanding
    \item Student runs code and observes results
    \item Errors are diagnosed and corrected with explanation
    \item Understanding is verified through physics questions
\end{itemize}

\subsection{Key Concepts Learned}

\begin{enumerate}
    \item \textbf{Python OOP}: Classes, methods, \texttt{self}, private methods
    (\texttt{\_validate}), DRY principle
    \item \textbf{Numerical methods}: Finite difference, sparse matrices,
    power iteration, convergence criteria
    \item \textbf{Nuclear physics}: 2-group diffusion, multi-group energy
    structure, boundary conditions, k-eff, flux distributions
    \item \textbf{Software engineering}: Git version control, unit testing
    with pytest, YAML input files, modular architecture
    \item \textbf{Nuclear data}: ENDF/B libraries, HDF5 format, group
    condensation, flux-weighted averages, OpenMC integration
    \item \textbf{Benchmarking}: IAEA 2D PWR benchmark, error analysis,
    pcm units, mesh convergence studies
\end{enumerate}

\subsection{Student Preferences Noted}

\begin{itemize}
    \item Prefers \textbf{fill-in-the-blank} exercises over full code
    \item Learns best from \textbf{visual explanations} (tables, diagrams)
    \item Strong \textbf{physics intuition} --- uses physical reasoning to
    catch coding errors
    \item Runs code from \textbf{ANGKOR root directory} via terminal
    \item Uses \textbf{VS Code} with Python extension on Windows
    \item Python environment: \textbf{Anaconda / venv} on Windows
    \item Has \textbf{MCNP experience} --- concepts transfer well
    \item Has \textbf{OpenMC in Docker} (openmc\_with\_jupyter image)
\end{itemize}

% =============================================================================
\section{Development Roadmap}
% =============================================================================

\subsection{Phase 1 --- Lattice Physics Code (Current)}

\begin{itemize}[leftmargin=*]
    \item[[x]] 1D slab finite difference solver
    \item[[x]] 2D rectangular geometry engine
    \item[[x]] 2-group and G-group neutron diffusion
    \item[[x]] IAEA 2D PWR quarter-core benchmark
    \item[[x]] YAML input system
    \item[[x]] Professional output module
    \item[[x]] Unit test suite
    \item[[TODO]] CMFD acceleration (in progress)
    \item[[TODO]] Cylindrical fuel pin geometry
    \item[[TODO]] 3D Cartesian geometry
\end{itemize}

\subsection{Phase 2 --- Neutronics Accuracy}

\begin{itemize}[leftmargin=*]
    \item[[TODO]] Transport correction for diffusion coefficients (B1/P1)
    \item[[TODO]] Adjoint flux solver (reactivity worth, perturbation theory)
    \item[[TODO]] Proper assembly XS generation (full assembly + B$^2$ search)
    \item[[TODO]] SP3 approximation for material interfaces
    \item[[TODO]] OECD/NEA C5G7 benchmark (7-group transport reference)
\end{itemize}

\subsection{Phase 3 --- Core Simulation}

\begin{itemize}[leftmargin=*]
    \item[[TODO]] Nodal Expansion Method (NEM/ANM)
    \item[[TODO]] Burnup and isotopic depletion (Bateman equations)
    \item[[TODO]] Reactivity feedback (Doppler, moderator density, boron)
    \item[[TODO]] Control rod worth and shutdown margin
    \item[[TODO]] Full PWR/BWR/SMR core models
\end{itemize}

\subsection{Phase 4 --- Advanced Features}

\begin{itemize}[leftmargin=*]
    \item[[TODO]] Full ENDF/B cross section processing pipeline
    \item[[TODO]] Thermal hydraulics coupling (single channel)
    \item[[TODO]] Verification against MCNP and Serpent
    \item[[TODO]] ML surrogate for cross section parameterization
    \item[[TODO]] Publication in nuclear engineering journal
\end{itemize}

% =============================================================================
\section{Key Technical References}
% =============================================================================

\begin{enumerate}
    \item IAEA, ``Benchmark Problem Book,'' ANL-7416, Argonne National
    Laboratory, 1977.

    \item J.J. Duderstadt and L.J. Hamilton, \textit{Nuclear Reactor Analysis},
    John Wiley \& Sons, 1976.

    \item K.S. Smith, ``Nodal Method Storage Reduction by Non-Linear
    Iteration,'' \textit{Trans. Am. Nucl. Soc.}, 44, 265, 1983.
    (CMFD foundation)

    \item Romano et al., ``OpenMC: A State-of-the-Art Monte Carlo Code for
    Research and Development,'' \textit{Ann. Nucl. Energy}, 82, 90-97, 2015.

    \item Stamm'ler and Abbate, \textit{Methods of Steady-State Reactor
    Physics in Nuclear Design}, Academic Press, 1983.
\end{enumerate}

% =============================================================================
\section{Quick Reference --- Running ANGKOR}
% =============================================================================

\subsection{Installation}

\begin{lstlisting}[language=bash, style=pythonstyle]
# Clone repository
git clone https://github.com/muthboravy007/ANGKOR.git
cd ANGKOR

# Create virtual environment
python -m venv venv
venv\Scripts\activate   # Windows

# Install dependencies
pip install -r requirements.txt
\end{lstlisting}

\subsection{Running the IAEA Benchmark}

\begin{lstlisting}[language=bash, style=pythonstyle]
python main.py input/iaea_2d.yaml
# Expected: k-eff = 0.998366, runtime ~43s
\end{lstlisting}

\subsection{Running Tests}

\begin{lstlisting}[language=bash, style=pythonstyle]
pytest tests/ -v
# Expected: 5/5 tests passing
\end{lstlisting}

\subsection{Git Workflow}

\begin{lstlisting}[language=bash, style=pythonstyle]
git pull origin main   # Always pull first!
git add <files>
git commit -m "description"
git push
\end{lstlisting}

% =============================================================================
\section{Glossary}
% =============================================================================}

\begin{center}
\begin{longtable}{p{3.5cm} p{9.5cm}}
\toprule
\textbf{Term} & \textbf{Definition} \\
\midrule
k-eff & Effective multiplication factor. Ratio of neutrons produced to neutrons
lost per generation. $k=1$: critical; $k>1$: supercritical; $k<1$: subcritical. \\[4pt]
pcm & Per cent mille. Unit for k-eff differences: $1\text{ pcm} = 10^{-5}$. \\[4pt]
Diffusion coefficient & $D = 1/(3\Sigma_{tr})$. Controls how far neutrons
diffuse. Larger D = more diffusion. \\[4pt]
5-point stencil & The finite difference approximation connecting each cell to
its 4 neighbors (N, S, E, W). \\[4pt]
Power iteration & Iterative method to find the dominant eigenvalue (k-eff) and
eigenvector (flux) of the diffusion equation. \\[4pt]
Dominance ratio & Ratio of first to fundamental eigenvalue. DR close to 1
means slow convergence. \\[4pt]
CMFD & Coarse Mesh Finite Difference. Acceleration technique using a coarse
mesh to speed up fine mesh convergence. \\[4pt]
d-hat & CMFD correction coefficient ensuring coarse mesh currents match fine
mesh currents at interfaces. \\[4pt]
Group condensation & Collapsing continuous-energy cross sections to few-group
constants using flux-weighted averages. \\[4pt]
Two-step method & Industry standard: (1) lattice physics generates few-group
XS, (2) core diffusion uses those XS. \\[4pt]
1/v law & Cross sections proportional to $1/v$ (inverse velocity) at thermal
energies. Thermal neutrons have much higher fission XS than fast neutrons. \\
\bottomrule
\end{longtable}
\end{center}

% =============================================================================
\vfill
\begin{center}
\rule{0.5\linewidth}{0.5pt}\\[0.5cm]
{\small ANGKOR --- Advanced Neutron Group-diffusion K-eigenvalue Of Reactors}\\
{\small Institute of Technology of Cambodia (ITC) $\cdot$ Phnom Penh, Cambodia}\\
{\small \url{https://github.com/muthboravy007/ANGKOR}}\\[0.3cm]
{\small\textit{``Named after Angkor Wat --- because Cambodia's scientific legacy continues.''}}
\end{center}

\end{document}
